\begin{textsample}{POS Dim 1 – 2013 – Score 23.00 – 2013/t002142.txt}  \label{ex:f1_pos_042}
Now, extinction is a different kind of death.
It ' s bigger.
We didn ' t really realize that until 1914, when \textbf{the} last passenger pigeon, a female named Martha, died at \textbf{the} Cincinnati zoo.
This had been \textbf{the} most abundant bird in \textbf{the} world that ' d been in North America for six million years.
Suddenly it wasn ' t here at all.
Flocks that were a mile wide and 400 miles long used to darken \textbf{the} sun.
Aldo Leopold said this was a biological storm, a feathered tempest.
And indeed it was a keystone species that enriched \textbf{the} entire \textbf{eastern} deciduous forest, from \textbf{the} Mississippi to \textbf{the} Atlantic, from Canada down to \textbf{the} Gulf.
But it went from five billion birds to zero in just a couple decades.
What happened?
Well, commercial hunting happened.
These birds were hunted for meat that was sold by \textbf{the} ton, and it was easy to do because when those big flocks came down to \textbf{the} ground, they were so dense that hundreds of hunters and netters could show up and slaughter them by \textbf{the} tens of thousands.
It was \textbf{the} cheapest source of protein in America.
By \textbf{the} end of \textbf{the} century, there was nothing left but these beautiful skins in museum specimen drawers.
There ' s an upside to \textbf{the} story.
This made people realize that \textbf{the} same thing was about to happen to \textbf{the} American bison, and so these birds saved \textbf{the} buffalos.
But a lot of other animals weren ' t saved. \textbf{The} Carolina parakeet was a parrot that lit up backyards everywhere.
It was hunted to death for its feathers.
There was a bird that people liked on \textbf{the} East Coast called \textbf{the} heath hen.
It was loved.
They tried to protect it.
It died anyway.
A local newspaper spelled out,"There is no survivor, there is no future, there is no life to be recreated in this form ever again."There ' s a sense of deep tragedy that goes with these things, and it happened to lots of birds that people loved.
It happened to lots of mammals.
Another keystone species is a \textbf{famous} animal called \textbf{the} European aurochs.
There was sort of a movie made about it recently.
And \textbf{the} aurochs was like \textbf{the} bison.
This was an animal that basically kept \textbf{the} forest mixed with grasslands across \textbf{the} entire Europe and Asian \textbf{continent}, from Spain to Korea. \textbf{The} documentation of this animal goes back to \textbf{the} Lascaux cave paintings. \textbf{The} extinctions still go on.
There ' s an ibex in Spain called \textbf{the} bucardo.
It went extinct in 2000.
There was a marvelous animal, a marsupial wolf called \textbf{the} thylacine in Tasmania, south of Australia, called \textbf{the} Tasmanian \textbf{tiger}.
It was hunted until there were just a few \textbf{left} to die in zoos.
A little bit of film was shot.
Sorrow, anger, mourning.
Don ' t mourn.
Organize.
What if you could find out that, using \textbf{the} DNA in museum specimens, \textbf{fossils} maybe up to 200, 000 years old could be used to bring species back, what would you do?
Where would you start?
Well, you ' d start by finding out if \textbf{the} biotech is really there.
I started with my wife, Ryan Phelan, who ran a biotech business called DNA Direct, and through her, one of her colleagues, George Church, one of \textbf{the} leading genetic engineers who turned out to be also obsessed with passenger pigeons and a lot of confidence that methodologies he was working on might actually do \textbf{the} deed.
So he and Ryan organized and hosted a meeting at \textbf{the} Wyss Institute in Harvard bringing together specialists on passenger pigeons, conservation ornithologists, bioethicists, and fortunately passenger pigeon DNA had already been sequenced by a molecular biologist named Beth Shapiro.
All she needed from those specimens at \textbf{the} Smithsonian was a little bit of toe pad tissue, because down in there is what is called ancient DNA.
It ' s DNA which is pretty badly fragmented, but with good \textbf{techniques} now, you can basically reassemble \textbf{the} whole \textbf{genome}.
Then \textbf{the} question is, can you reassemble, with that \textbf{genome}, \textbf{the} whole bird?
George Church thinks you can.
So in his book,"Regenesis,"which I recommend, he has a chapter on \textbf{the} science of bringing back extinct species, and he has a machine called \textbf{the} Multiplex Automated \textbf{Genome} Engineering machine.
It ' s kind of like an \textbf{evolution} machine.
You try combinations of genes that you write at \textbf{the} cell level and then in organs on a chip, and \textbf{the} ones that win, that you can then put into a living organism.
It ' ll work. \textbf{The} precision of this, one of George ' s \textbf{famous} unreadable \textbf{slides}, nevertheless points out that there ' s a level of precision here right down to \textbf{the} individual base pair. \textbf{The} passenger pigeon has 1. 3 billion base pairs in its \textbf{genome}.
So what you ' re getting is \textbf{the} capability now of replacing one gene with another variation of that gene.
It ' s called an allele.
Well that ' s what happens in normal hybridization anyway.
So this is a form of synthetic hybridization of \textbf{the} \textbf{genome} of an extinct species with \textbf{the} \textbf{genome} of its closest living \textbf{relative}.
Now along \textbf{the} way, George points out that his technology, \textbf{the} technology of synthetic \textbf{biology}, is currently accelerating at four times \textbf{the} rate of Moore ' s Law.
It ' s been doing that since 2005, and it ' s likely to continue. \textbf{Okay}, \textbf{the} closest living \textbf{relative} of \textbf{the} passenger pigeon is \textbf{the} band-tailed pigeon.
They ' re abundant.
There ' s some around here.
Genetically, \textbf{the} band- tailed pigeon already is mostly living passenger pigeon.
There ' s just some bits that are band-tailed pigeon.
If you replace those bits with passenger pigeon bits, you ' ve got \textbf{the} extinct bird back, cooing at you.
Now, there ' s work to do.
You have to figure out exactly what genes matter.
So there ' s genes for \textbf{the} short tail in \textbf{the} band-tailed pigeon, genes for \textbf{the} long tail in \textbf{the} passenger pigeon, and so on with \textbf{the} red eye, peach-colored breast, flocking, and so on.
Add them all up and \textbf{the} result won ' t be perfect.
But it should be be perfect enough, because nature doesn ' t do perfect either.
So this meeting in Boston led to three things.
First off, Ryan and I decided to create a nonprofit called Revive and Restore that would push de-extinction generally and try to have it go in a responsible way, and we would push ahead with \textbf{the} passenger pigeon.
Another direct result was a young grad student named Ben Novak, who had been obsessed with passenger pigeons since he was 14 and had also learned how to work with ancient DNA, himself sequenced \textbf{the} passenger pigeon, using money from his family and friends.
We hired him full-time.
Now, this photograph I took of him last year at \textbf{the} Smithsonian, he ' s looking down at Martha, \textbf{the} last passenger pigeon alive.
So if he ' s successful, she won ' t be \textbf{the} last. \textbf{The} third result of \textbf{the} Boston meeting was \textbf{the} realization that there are scientists all over \textbf{the} world working on various forms of de-extinction, but they ' d never met each other.
And National Geographic got interested because National Geographic has \textbf{the} theory that \textbf{the} last century, discovery was basically finding things, and in this century, discovery is basically making things.
De-extinction falls in that category.
So they hosted and funded this meeting.
And 35 scientists, they were conservation biologists and molecular biologists, basically meeting to see if they had work to do together.
Some of these conservation biologists are pretty radical.
There ' s three of them who are not just re-creating ancient species, they ' re recreating extinct ecosystems in northern Siberia, in \textbf{the} Netherlands, and in Hawaii.
Henri, from \textbf{the} Netherlands, with a Dutch last name I won ' t try to pronounce, is working on \textbf{the} aurochs. \textbf{The} aurochs is \textbf{the} ancestor of all domestic cattle, and so basically its \textbf{genome} is alive, it ' s just unevenly distributed.
So what they ' re doing is working with seven breeds of primitive, hardy-looking cattle like that Maremmana primitivo on \textbf{the} top there to rebuild, over time, with selective back-breeding, \textbf{the} aurochs.
Now, re-wilding is moving faster in Korea than it is in America, and so \textbf{the} plan is, with these re-wilded areas all over Europe, they will introduce \textbf{the} aurochs to do its old job, its old ecological role, of clearing \textbf{the} somewhat barren, closed-canopy forest so that it has these biodiverse meadows in it.
Another amazing story came from Alberto Fern{\EmojiFont �}{\EmojiFont �}ndez-Arias.
Alberto worked with \textbf{the} bucardo in Spain. \textbf{The} last bucardo was a female named Celia who was still alive, but then they captured her, they got a little bit of tissue from her ear, they cryopreserved it in liquid nitrogen, released her back into \textbf{the} wild, but a few months later, she was found dead under a fallen tree.
They took \textbf{the} DNA from that ear, they planted it as a \textbf{cloned} egg in a goat, \textbf{the} pregnancy came to term, and a live baby bucardo was born.
It was \textbf{the} first de-extinction in history.
It was short-lived.
Sometimes interspecies clones have respiration problems.
This one had a malformed lung and died after 10 minutes, but Alberto was confident that \textbf{cloning} has moved along well since then, and this will move ahead, and eventually there will be a population of bucardos back in \textbf{the} mountains in northern Spain.
Cryopreservation pioneer of great depth is Oliver Ryder.
At \textbf{the} San Diego zoo, his frozen zoo has collected \textbf{the} tissues from over 1, 000 species over \textbf{the} last 35 years.
Now, when it ' s frozen that deep, minus 196 degrees Celsius, \textbf{the} cells are intact and \textbf{the} DNA is intact.
They ' re basically viable cells, so someone like Bob Lanza at Advanced Cell Technology took some of that tissue from an endangered animal called \textbf{the} Javan banteng, put it in a cow, \textbf{the} cow went to term, and what was born was a live, healthy baby Javan banteng, who thrived and is still alive. \textbf{The} most exciting thing for Bob Lanza is \textbf{the} ability now to take any kind of cell with induced pluripotent stem cells and turn it into germ cells, like sperm and eggs.
So now we go to Mike McGrew who is a scientist at Roslin Institute in Scotland, and Mike ' s doing miracles with birds.
So he ' ll take, say, falcon skin cells, fibroblast, turn it into induced pluripotent stem cells.
Since it ' s so pluripotent, it can become germ plasm.
He then has a way to put \textbf{the} germ plasm into \textbf{the} embryo of a chicken egg so that that chicken will have, basically, \textbf{the} gonads of a falcon.
You get a male and a female each of those, and out of them comes falcons.
Real falcons out of slightly doctored chickens.
Ben Novak was \textbf{the} youngest scientist at \textbf{the} meeting.
He showed how all of this can be put together. \textbf{The} sequence of events : he ' ll put together \textbf{the} \textbf{genomes} of \textbf{the} band-tailed pigeon and \textbf{the} passenger pigeon, he ' ll take \textbf{the} \textbf{techniques} of George Church and get passenger pigeon DNA, \textbf{the} \textbf{techniques} of Robert Lanza and Michael McGrew, get that DNA into chicken gonads, and out of \textbf{the} chicken gonads get passenger pigeon eggs, squabs, and now you ' re getting a population of passenger pigeons.
It does raise \textbf{the} question of, they ' re not going to have passenger pigeon parents to teach them how to be a passenger pigeon.
So what do you do about that?
Well birds are pretty hard-wired, as it happens, so most of that is already in their DNA, but to \textbf{supplement} it, part of Ben ' s idea is to use homing pigeons to help train \textbf{the} young passenger pigeons how to flock and how to find their way to their old nesting grounds and \textbf{feeding} grounds.
There were some conservationists, really \textbf{famous} conservationists like Stanley Temple, who is one of \textbf{the} founders of conservation \textbf{biology}, and Kate Jones from \textbf{the} IUCN, which does \textbf{the} Red List.
They ' re excited about all this, but they ' re also concerned that it might be competitive with \textbf{the} extremely important efforts to protect endangered species that are still alive, that haven ' t gone extinct yet.
You see, you want to work on protecting \textbf{the} animals out there.
You want to work on getting \textbf{the} market for ivory in Asia down so you ' re not using 25, 000 \textbf{elephants} a year.
But at \textbf{the} same time, conservation biologists are realizing that bad news bums people out.
And so \textbf{the} Red List is really important, keep track of what ' s endangered and critically endangered, and so on.
But they ' re about to create what they call a Green List, and \textbf{the} Green List will have species that are doing fine, thank you, species that were endangered, like \textbf{the} bald eagle, but they ' re much better off now, thanks to everybody ' s good work, and protected areas around \textbf{the} world that are very, very well managed.
So basically, they ' re \textbf{learning} how to build on good news.
And they see reviving extinct species as \textbf{the} kind of good news you might be able to build on.
Here ' s a couple related examples.
Captive breeding will be a major part of bringing back these species. \textbf{The} California condor was down to 22 birds in 1987.
Everybody thought is was finished.
Thanks to captive breeding at \textbf{the} San Diego Zoo, there ' s 405 of them now, 226 are out in \textbf{the} wild.
That technology will be used on de-extincted animals.
Another success story is \textbf{the} mountain gorilla in Central Africa.
In 1981, Dian Fossey was sure they were going extinct.
There were just 254 \textbf{left}.
Now there are 880.
They ' re increasing in population by three percent a year. \textbf{The} secret is, they have an eco-tourism program, which is absolutely brilliant.
So this photograph was taken last month by Ryan with an iPhone.
That ' s how comfortable these wild gorillas are with visitors.
Another interesting project, though it ' s going to need some help, is \textbf{the} northern white rhinoceros.
There ' s no breeding pairs left.
But this is \textbf{the} kind of thing that a wide variety of DNA for this animal is available in \textbf{the} frozen zoo.
A bit of \textbf{cloning}, you can get them back.
So where do we go from here?
These have been private meetings so far.
I think it ' s time for \textbf{the} subject to go public.
What do people think about it?
You know, do you want extinct species back?
Do you want extinct species back?
Tinker Bell is going to come fluttering down.
It is a Tinker Bell moment, because what are people excited about with this?
What are they concerned about?
We ' re also going to push ahead with \textbf{the} passenger pigeon.
So Ben Novak, even as we speak, is joining \textbf{the} group that Beth Shapiro has at UC Santa Cruz.
They ' re going to work on \textbf{the} \textbf{genomes} of \textbf{the} passenger pigeon and \textbf{the} band-tailed pigeon.
As that data matures, they ' ll send it to George Church, who will work his magic, get passenger pigeon DNA out of that.
We ' ll get help from Bob Lanza and Mike McGrew to get that into germ plasm that can go into chickens that can produce passenger pigeon squabs that can be raised by band-tailed pigeon parents, and then from then on, it ' s passenger pigeons all \textbf{the} way, maybe for \textbf{the} next six million years.
You can do \textbf{the} same thing, as \textbf{the} costs come down, for \textbf{the} Carolina parakeet, for \textbf{the} great auk, for \textbf{the} heath hen, for \textbf{the} ivory- billed woodpecker, for \textbf{the} Eskimo curlew, for \textbf{the} Caribbean \textbf{monk} \textbf{seal}, for \textbf{the} woolly mammoth.
Because \textbf{the} fact is, humans have made a huge hole in nature in \textbf{the} last 10, 000 years.
We have \textbf{the} ability now, and maybe \textbf{the} moral obligation, to repair some of \textbf{the} damage.
Most of that we ' ll do by expanding and protecting wildlands, by expanding and protecting \textbf{the} populations of endangered species.
But some species that we killed off totally we could consider bringing back to a world that misses them.
Thank you.
Chris Anderson : Thank you.
I ' ve got a question.
So, this is an emotional topic.
Some people stand.
I suspect there are some people out there sitting, kind of asking tormented questions, almost, about, well, wait, wait, wait, wait, wait, wait a minute, there ' s something wrong with mankind interfering in nature in this way.
There ' s going to be unintended consequences.
You ' re going to uncork some sort of Pandora ' s box of who-knows-what.
Do they have a point?
Stewart Brand : Well, \textbf{the} earlier point is we interfered in a big way by making these animals go extinct, and many of them were keystone species, and we changed \textbf{the} whole ecosystem they were in by letting them go.
Now, there ' s \textbf{the} shifting baseline problem, which is, so when these things come back, they might replace some birds that are there that people really know and love.
I think that ' s, you know, part of how it ' ll work.
This is a long, slow process — One of \textbf{the} things I like about it, it ' s multi-generation.
We will get woolly mammoths back.
CA : Well it feels like both \textbf{the} conversation and \textbf{the} potential here are pretty thrilling.
Thank you so much for presenting.
SB : Thank you.
CA : Thank you.

% matched lemmas: biology, clon, continent, eastern, elephant, evolution, famous, feeding, fossil, genome, learning, left, monk, okay, relative, seal, slide, supplement, technique, the, tiger
\end{textsample}
